% Context from chapters/approximation.tex; for mathematical reference, not compilation.
\section{Wavelet Approximation of Piecewise Smooth Functions}

Wavelets improve approximation near singularities because their support is localized.

                                        
\subsection{Decay of Wavelet Coefficients}
\label{subsect-wavdecay}

To approximate smooth regions efficiently, choose wavelets with sufficiently many vanishing moments, denoted by $p$:
\eq{
	\foralls k=0,\ldots,p-1, \; \int \psi(x) x^k \d x = 0.
}
Choose $p\geq\alpha$, where $\alpha$ is the regularity of the signal outside singularities (for instance jumps or kinks).

\myfigure{
\image{approximation}{.55}{singular-part-1d}
\image{approximation}{.25}{singular-part-2d}
}{ 
	Singular and regular regions of a signal (left) and an image (right). 
}{fig-singular-part}

To quantify the approximation error decay for piecewise smooth signals \eqref{eq-model-piece-smooth} and images \eqref{eq-model-piece-smooth-2d}, one treats wavelets supported in smooth regions separately from those crossing singularities. Figure \ref{fig-singular-part} locates the singular and regular regions of a signal and an image.

\begin{prop}
Assume a compactly supported wavelet with $p\geq\alpha$ vanishing moments (all moments of total degree less than $p$ in dimension $d$). If $f$ admits a $C^\alpha$ extension to a neighborhood of the support, one has
\eql{\label{eq-wavcoef-decay-1}
	|\dotp{f}{\psi_{j,n}}| \leq C_f2^{j(\alpha+d/2)}\int|t|^\alpha|\psi(t)|\,\d t.
}
Without the smoothness assumption, every bounded $f$ satisfies
\eql{\label{eq-wavcoef-decay-2}
	|\dotp{f}{\psi_{j,n}}| \leq \norm{f}_\infty \norm{\psi}_1 2^{j \frac{d}{2}}.
}
\end{prop}

\begin{proof}
Under these assumptions, one can perform a Taylor expansion of $f$ around the point $2^j n$
\eq{
	f(x) = P(x-2^j n) + R(x-2^jn) = P(2^j t) + R(2^j t)
}
where $\deg(P)<\al$ and 
\eq{
	|R(x)|\leq C_f \norm{x}^\al.
}
One then bounds the wavelet coefficient
\eq{
	\dotp{f}{\psi_{j,n}} = \frac{1}{2^{j \frac{d}{2}}} \int f(x) \psi\pa{ \frac{x-2^j n}{2^j} } \d x 
		= 2^{j \frac{d}{2}} \int R(2^j t) \psi(t) \d t
}
where we substituted $t = \frac{x-2^j n}{2^j}$. This proves~\eqref{eq-wavcoef-decay-1}. The bound~\eqref{eq-wavcoef-decay-2} follows directly by bounding the function in the inner product.
\end{proof}

                                        
\subsection{1-D Piecewise Smooth Approximation}

For a piecewise smooth 1-D signal in the model~\eqref{eq-model-piece-smooth}, the large coefficients $\dotp{f}{\psi_{j,n}}$ concentrate near singularities. Denote the finite set of singular points by $\Ss \subset [0,1]$.


\begin{thm}\label{thm-approx-piecesmooth-1d}
For $\alpha>0$, let $f$ belong to the piecewise $C^\alpha$ model~\eqref{eq-model-piece-smooth}. With compactly supported orthonormal wavelets having at least $\alpha$ vanishing moments (and a compatible boundary construction), the best $M$-term error satisfies
\eql{\label{eq-decay-piecewav}
	\norm{f-f_M}^2 = O(M^{-2\al}).
}
\end{thm}

\begin{proof}
The proof is split into several parts.

  
\paragraph{Step 1. Separating regular and singular coefficients.}

At scale $2^j$, define the singular index set by the wavelets whose support crosses a singularity:
\eql{\label{eq-singsupport-1d}
	\Cc_j = \enscond{n}{ \supp(\psi_{j,n}) \cap \Ss \neq \emptyset }
}
Compact support and dyadic spacing imply a bound on its cardinality that is uniform in $j$:
\eq{
	|\Cc_j| \leq K |\Ss| = \text{constant}.
}
Using \eqref{eq-wavcoef-decay-1} for $d=1$, the decay of regular coefficients is bounded as
\eq{
	\foralls n \in \Cc_j^c, \quad \abs{\dotp{f}{\psi_{j,n}}} \leq  C2^{j(\alpha+1/2)}.
}
Using \eqref{eq-wavcoef-decay-2} for $d=1$, the decay of singular coefficients is bounded as
\eq{
	\foralls n \in \Cc_j, \quad \abs{\dotp{f}{\psi_{j,n}}} \leq  C 2^{j/2}.
}
For a sufficiently small threshold $T$, define two cutoff scales, one for regular coefficients and one for singular coefficients, as functions of $T$:
\eq{
	2^{j_1} = (T/C)^{\frac{1}{\al+1/2}}
	\qandq
	2^{j_2} = (T/C)^{2}.
}
Round these cutoff scales to integers; this changes constants only. All subsequent scale sums are bounded above by the fixed coarsest scale $j_0=0$, and the finitely many scaling coefficients are retained. Figure~\ref{fig-singular-support-1d} shows a schematic segmentation of the set of wavelet coefficients into regular and singular parts, together with the cut-off scales.

\myfigure{
\image{approximation}{.6}{singular-support-1d}
\image{approximation}{.3}{singular-support-legend}
}{ 
	Segmentation of the wavelet coefficients into regular and singular parts.  
}{fig-singular-support-1d}

  
\paragraph{Step 2. Bounding the error.}
These cut-off scales allow us to define a comparison approximation
\eql{\label{eq-hand-made-approx}
	\tilde f_M=P_{V_0}f+\sum_{j_2\leq j\leq0} \sum_{n \in \Cc_j} \dotp{f}{\psi_{j,n}} \psi_{j,n} +  \sum_{j_1\leq j\leq0} \sum_{n \in \Cc_j^c} \dotp{f}{\psi_{j,n}} \psi_{j,n}.
}
The error of the comparison $M$-term approximation $\tilde f_M$ bounds the best $M$-term error from above:
\eqml{\label{eq-wavapprox-1d-1}
	\norm{f-f_M}^2 &\leq \norm{f-\tilde f_M}^2 \leq \sum_{j<j_2, n \in \Cc_j} |\dotp{f}{\psi_{j,n}}|^2 + 
\sum_{j<j_1, n \in \Cc_j^c} |\dotp{f}{\psi_{j,n}}|^2 \\
	& \leq \sum_{j < j_2} (K |\Ss|) \times C^2 2^{j} + \sum_{j < j_1} 2^{-j} \times  C^2 2^{j(2\al+1)} \\
	& = O( 2^{j_2} + 2^{2\al j_1} ) = O( T^2 + T^{ \frac{2\al}{\al+1/2} } ) = O( T^{ \frac{2\al}{\al+1/2} } ).
}

  
\paragraph{Step 3. Counting the Coefficients.}
The number of coefficients needed to build the approximating signal $\tilde f_M$ is
\eqml{\label{eq-wavapprox-1d-2}
	M &\leq1+\sum_{j_2\leq j\leq0} |\Cc_j| +
\sum_{j_1\leq j\leq0} |\Cc_j^c| \leq
1+\sum_{j_2\leq j\leq0} K |\Ss| +
\sum_{j_1\leq j\leq0} 2^{-j} \\
	&= O( |\log(T)| + T^{\frac{-1}{\al+1/2}} ) 
=  O(T^{\frac{-1}{\al+1/2}} ).
}
  
\paragraph{Step 4. Putting everything together.}
Given a coefficient budget $M$, choose $T$ to be a sufficiently large constant times $M^{-(\alpha+1/2)}$. The coefficient-count bound then uses at most $M$ terms, and the error bound is $O(M^{-2\alpha})$.
\end{proof}

For piecewise Lipschitz signals, the wavelet rate is faster than the jump-limited $O(M^{-1})$ Fourier bound in~\eqref{eq-fourier-piecewisesmooth}. It also matches the smooth-signal rate~\eqref{eq-sovolev-signal-decay}: finitely many singularities do not worsen the asymptotic exponent in one dimension. The rate~\eqref{eq-decay-piecewav} is asymptotically optimal for this class.

\myfigure{
\image{approximation}{.6}{wav-approx-1d-1}\\
\image{approximation}{.6}{wav-approx-1d-2}\\
\image{approximation}{.6}{w